\documentclass[11pt, a4paper]{article}
\usepackage[margin=2.7cm]{geometry}
\usepackage{amsmath, amssymb, amsthm}
\usepackage{booktabs}
\usepackage{microtype}
\usepackage{xcolor}
\usepackage[colorlinks=true,linkcolor=blue!60!black,citecolor=blue!60!black,urlcolor=blue!60!black]{hyperref}
\hypersetup{pdftitle={Identification of the order distribution of linear
systems by line-spectrum estimation on logarithmic spirals}}

\newtheorem{theorem}{Theorem}[section]
\newtheorem{proposition}[theorem]{Proposition}
\newtheorem{lemma}[theorem]{Lemma}
\newtheorem{corollary}[theorem]{Corollary}
\theoremstyle{definition}
\newtheorem{definition}[theorem]{Definition}
\newtheorem{example}[theorem]{Example}
\theoremstyle{remark}
\newtheorem{remark}[theorem]{Remark}

\newcommand{\R}{\mathbb{R}}
\newcommand{\Z}{\mathbb{Z}}
\newcommand{\C}{\mathbb{C}}
\newcommand{\dd}{\,\mathrm{d}}
\newcommand{\ii}{\mathrm{i}}
\newcommand{\ee}{\mathrm{e}}

\title{Identification of the order distribution of linear systems\\
by line-spectrum estimation on logarithmic spirals}
\author{Ing.\ Ulises Merlan\thanks{Ingeniero Electr\'onico, Universidad
Tecnol\'ogica Nacional, Facultad Regional Buenos Aires. Email:
\texttt{ulimerlan@frba.utn.edu.ar}.}}
\date{July 2026}

\begin{document}
\maketitle

\begin{abstract}
A linear time-invariant system whose transfer function admits the
distributed-order representation $H(s)=\int_\R \mathcal{A}(\alpha)\,
s^{\alpha} \dd\alpha$ is characterized by its \emph{order distribution}
$\mathcal{A}$: integer-order systems correspond to Dirac combs supported on
$\Z$ whose weights are Laurent--Taylor coefficients, while fractional-order
systems place mass at non-integer $\alpha$s. Because the order integral is a
Laplace transform of $\mathcal{A}$ in the variable $\log s$, sampling $H$ on
the unit circle $s=\ee^{\ii u}$ and applying an FFT suggests itself as a
direct, closed-form estimator of $\mathcal{A}$ from measured data. We show
that this natural estimator is subject to two structural obstructions:
(i)~a truncation bias, computed here in closed form, which corrupts every
coefficient whenever the measured response has a non-zero steady state; and
(ii)~an aliasing theorem stating that any transform computed from time-domain
data via $s=\ee^{\ii u}$ is exactly $2\pi$-periodic, so its inverse Fourier
transform is supported on $\Z$ and non-integer orders can \emph{never} appear
as spectral lines, only as slowly decaying integer combs. We then repair
both defects: subtracting the measured steady state reduces the integer-order
problem to classical Cauchy-integral coefficient extraction
(Lyness--Moler), and replacing the circle by a \emph{logarithmic spiral}
$s=\ee^{(\sigma+\ii)u}$, $\sigma\neq0$, breaks the periodicity, turning
order recovery into a line-spectrum estimation problem solvable by
matrix-pencil methods. Numerical experiments recover the seven leading
Taylor coefficients of a second-order test system from its sampled step
response with error ${\sim}10^{-8}$, and recover the order pair
$(-1,-\tfrac12)$ (one integer, one fractional line) of a signal, from raw
time samples with error ${\sim}2\times10^{-9}$ (noise-free) and
${\sim}5\times10^{-5}$ under $10^{-4}$ relative measurement noise. As a
calibration benchmark, the estimator recovers the complex dimensions of
the Cantor fractal string---known in closed form---from geometric
counting data alone, with damping estimates matching the exact value
$D=\log2/\log3$ to six decimals, and, from spectral counting data,
recovers the same dimensions together with channel gains matching
independently computed values of the Riemann zeta function to four
decimals. An open-source implementation is provided.
\end{abstract}

\noindent\textbf{Keywords:} fractional calculus, distributed-order systems,
order distribution, system identification, Mellin transform, matrix pencil,
Cauchy integrals, fractal strings, complex dimensions.

\noindent\textbf{2020 MSC:} 26A33, 93B30, 65F99.

\section{Introduction}\label{sec:intro}

Fractional-order models---systems governed by derivatives
$D^{\alpha}$ of non-integer order $\alpha$---arise in viscoelasticity,
electrochemistry, anomalous diffusion, and control
\cite{DingReview2021, JiaoChenPodlubny2012}. A strictly more general class
allows a \emph{continuum} of orders: distributed-order operators
$\int \mathcal{A}(\alpha) D^{\alpha}y \dd\alpha$, introduced by Caputo
\cite{Caputo1969,Caputo1995, Caputo2001} and developed by Bagley and Torvik
\cite{BagleyTorvik2000a,BagleyTorvik2000b} and by Lorenzo and Hartley
\cite{LorenzoHartley2002}. In the Laplace domain (zero initial conditions)
such a system has transfer function
\begin{equation}\label{eq:orderintegral}
H(s) \;=\; \int_{\R} \mathcal{A}(\alpha)\, s^{\alpha} \dd\alpha ,
\end{equation}
and the density $\mathcal{A}$, called the \emph{order distribution}
\cite{HartleyLorenzo1999,HartleyLorenzo2003}, is the natural object of
identification: an integer-order rational system corresponds to
$\mathcal{A}$ being a finite or geometric-weight Dirac comb on the integers,
a commensurate fractional system to a comb on $q\Z$ for some $q\in(0,1)$,
and genuinely distributed systems to continuous densities.

Identifying $\mathcal{A}$ from measured data is a recognized problem.
Published approaches discretize the order axis and fit the weights by least
squares to frequency-response data
\cite{HartleyLorenzo2003,NazarianHaeri2010,DuanLi2021}, estimate discrete
fractional orders by nonlinear optimization or instrumental variables
\cite{VictorMalti2013,VictorMalti2022}, read asymptotic Bode slopes, or, in
the PDE setting, exploit analyticity of solutions in inverse-problem
frameworks \cite{LiuSunYamamoto2021}. What makes a \emph{closed-form}
alternative tempting is a structural identity already noted in
\cite{HartleyLorenzo1999}: the right-hand side of \eqref{eq:orderintegral}
is a (bilateral) Laplace transform of $\mathcal{A}$ with transform variable
$-\log s$. Consequently, on the unit circle $s=\ee^{\ii u}$ the identity
reads $H(\ee^{\ii u})=\int\mathcal{A}(\alpha)\ee^{\ii u\alpha}\dd\alpha$: a
Fourier transform. One is led to the estimator
\begin{equation}\label{eq:naiveinversion}
\mathcal{A}(\alpha) \;=\; \frac{1}{2\pi}\int_{\R}
H(\ee^{\ii u}) \, \ee^{-\ii u\alpha}\dd u ,
\end{equation}
computable by an FFT from samples of $H$ on the circle, which in turn can be
obtained from a measured step or impulse response by numerically evaluating
its Laplace transform at $s=\ee^{\ii u_k}$. Orders present in the system
would then appear as spectral lines of $\mathcal{A}$, integer or not.

This paper analyzes that program and establishes exactly how much of it
survives contact with measured data. Our contributions are:

\begin{enumerate}
\item \textbf{A corrected integer-order pipeline with an exact bias
formula.} The naive estimator applied to a step response with steady state
$y_\infty\neq0$ suffers a truncation bias
$-y_\infty(-A)^{m+1}/(m+1)!$ at integer order $m$, where $A$ is the
measurement window (Proposition~\ref{prop:bias}); the bias grows
exponentially in $A$ across the bins (like $\ee^{A}$ up to algebraic
factors) and dwarfs the true coefficients. Subtracting the measured steady
state before transforming removes it entirely and reduces the method to
Cauchy-integral coefficient extraction in the sense of Lyness and Moler
\cite{LynessMoler1967} (Section~\ref{sec:circle}).

\item \textbf{An aliasing obstruction.} Any function of $u$ computed from
time samples through $s=\ee^{\ii u}$ is exactly $2\pi$-periodic; therefore
its inverse Fourier transform is supported on $\Z$
(Theorem~\ref{thm:aliasing}). Non-integer orders are structurally invisible
to the unit-circle readout of black-box data: a branch cut manifests itself
only as an integer comb with slow $\mathcal{O}(1/m)$ decay, for which we
give a closed-form example. This clarifies which claims of the closed-form
program hold for transfer functions known analytically (where the Riemann
surface of $s^{\alpha}$ can be unwound) but fail for data.

\item \textbf{A logarithmic-spiral estimator for fractional orders.}
Sampling the data transform on $s=\ee^{(\sigma+\ii)u}$ with $\sigma\neq0$
makes $\log s$ single-valued and injective along the contour; a finite
order comb then produces an exponential sum in the sample index, and the
orders and weights follow from a matrix-pencil line-spectrum estimate
\cite{HuaSarkar1990} (Section~\ref{sec:spiral}). The estimator works
directly on raw time samples, with the contour restricted to an arc in the
right half-plane to keep the data transform convergent.

\item \textbf{Numerical validation and open code.} On the step response of
$H(s)=1/(s^{2}+4s+5)$, the corrected circle pipeline recovers
$a_0,\dots,a_6$ with maximum error $1.2\times10^{-8}$, the steady state
being estimated from the data itself; the spiral estimator recovers the
order pair $(-1,-\tfrac12)$ of $y(t)=1/\sqrt{\pi t}+2$ from time samples
with error $1.8\times10^{-9}$, degrading to $4.7\times10^{-5}$ under
$10^{-4}$ relative noise (Section~\ref{sec:experiments}).

\item \textbf{A calibration benchmark from fractal geometry.} The estimator
is validated against complex dimensions of self-similar fractal strings, a
setting where the ``uniform damping'' question---conjectural for the
Riemann zeta function---is a \emph{theorem}: for the Cantor string, the
estimator recovers eight complex dimensions from geometric data with all
damping estimates equal to the exact value $D=\log2/\log3$ to six
decimals, and, from spectral counting data alone, recovers the same
dimensions together with channel gains matching independently computed
values of $|\zeta|$ to four decimals in every mode---an acoustic
measurement of the zeta function realizing the classical spectral
factorization $\zeta_\nu=\zeta_{\mathcal{L}}\cdot\zeta$ as a live transfer
function (Section~\ref{sec:e5}). A further experiment exhibits the
estimator's channel gain vanishing, as theory predicts, at a zeta zero.
\end{enumerate}

Section~\ref{sec:related} positions these results against the
distributed-order identification literature and against the numerical
analysis of Cauchy-integral coefficient extraction; proofs of the two
propositions are collected in Appendix~\ref{app:proofs}.

\section{The order-distribution representation}\label{sec:representation}

Throughout, $s^{\alpha}=\ee^{\alpha\log s}$ with a branch of $\log$ fixed by
the contour under discussion, and $\mathcal{A}$ is a tempered distribution
on $\R$; \eqref{eq:orderintegral} is read in the distributional sense.

\begin{definition}\label{def:orderdist}
Let $H$ be analytic on an open annulus
$\Omega=\{r_1<|s|<r_2\}$, $0\le r_1<r_2\le\infty$. A tempered distribution
$\mathcal{A}$ on $\R$ is an \emph{order distribution for $H$ on $\Omega$} if
\[
H(\ee^{w})=\big\langle \mathcal{A}(\alpha), \ee^{\alpha w}\big\rangle
\]
for all $w$ with $\log r_1<\operatorname{Re}w<\log r_2$.
\end{definition}

\begin{example}\label{ex:laurent}
If $H$ is analytic on an annulus containing $|s|=1$ and single-valued, its
Laurent expansion $H(s)=\sum_{n\in\Z}a_n s^{n}$ converges there with
geometrically decaying $a_n$, and
\[
\mathcal{A}=\sum_{n\in\Z}a_n\, \delta(\alpha-n)
\]
is an order distribution for $H$: the classical case. For instance,
$H(s)=s^{2}+k$ gives $\delta(\alpha-2)+k\, \delta(\alpha)$;
$H(s)=(s+1)/s$ gives $\delta(\alpha)+\delta(\alpha+1)$;
$H(s)=1/(s^{2}+2)$ gives $\sum_{n\ge0}(-1)^{n}2^{-(n+1)}
\delta(\alpha-2n)$ on $|s|<\sqrt2$. A Caputo-type fractional term
contributes mass at its order: $H(s)=s^{1/2}$ has
$\mathcal{A}=\delta(\alpha-\tfrac12)$ on the universal cover of
$\C\setminus\{0\}$.
\end{example}

\begin{remark}[Non-uniqueness]\label{rem:annulus}
The order distribution depends on the annulus. E.g.\
$1/(s^{2}+2)$ also equals $\sum_{n\ge0}(-1)^{n}2^{n}s^{-2n-2}$ on
$|s|>\sqrt2$, a different comb. Statements such as ``the positions of the
deltas reveal the orders of the system'' are therefore well-posed only once
a contour (equivalently, a strip of the Mellin variable) is fixed. In what
follows the annulus is always the one containing the sampling contour.
\end{remark}

Writing $s=\ee^{\ii u}$, $u\in\R$, Definition~\ref{def:orderdist} gives the
Fourier pair
\begin{equation}\label{eq:fourierpair}
H(\ee^{\ii u})=\int_\R \mathcal{A}(\alpha) \, \ee^{\ii u\alpha} \dd\alpha ,
\qquad
\mathcal{A}(\alpha)=\frac{1}{2\pi}\int_{\R} H(\ee^{\ii u})\,
\ee^{-\ii u\alpha}\dd u ,
\end{equation}
where for multivalued $H$ the left-hand side denotes the lift along the
universal cover (i.e.\ $u$ is \emph{not} reduced modulo $2\pi$). When
$\mathcal{A}=\sum_n a_n\, \delta(\alpha-n)$, the weights are given by the
Cauchy formula: substituting $s=\ee^{\ii u}$, $\dd s=\ii s\dd u$,
\begin{equation}\label{eq:cauchy}
a_n\;=\; \frac{1}{2\pi\ii}
\oint_{|s|=1} \frac{H(s)}{s^{n+1}}\dd s ,
\end{equation}
with the kernel $s^{-\alpha-1}$ (not $s^{-\alpha}$): the Jacobian
$\dd s/s$ is what aligns the continuous inversion with the classical
coefficient formula. Under $s=\ee^{-p}$, \eqref{eq:orderintegral} is a
bilateral Laplace transform of $\mathcal{A}$; the representation and its
inversion are thus an instance of the Mellin correspondence
\cite{LuchkoKiryakova2013}, and the interpolation of power-series
coefficients to continuous order is the territory of Ramanujan's master
theorem. None of this section is new mathematics; it fixes notation and
conventions for the data-driven question that follows.

\section{Recovery from measured data on the unit circle}\label{sec:circle}

Let the system be excited by a unit step, $X(s)=1/s$, and let
$y(t)$, $t\in[0,A]$, be the measured response, so that
$Y(s)=H(s)/s$. The data provide the \emph{truncated} transform
\begin{equation}\label{eq:truncated}
\widehat{Y}_{\!A}(s)\;=\;\int_0^{A} y(\tau)\,\ee^{-s\tau}\dd\tau ,
\end{equation}
an entire function of $s$ computable by quadrature at any $s_k=\ee^{\ii u_k}$.
The naive pipeline evaluates \eqref{eq:truncated} on the circle, multiplies
by $s$ to undo the input, and applies a DFT across $u$ to read off the comb
weights. Two elementary identities explain its behavior exactly.

\begin{proposition}[Aliased-moment identity]\label{prop:moments}
Let $u_k=2\pi k/N$, $k=0, \dots,N-1$, and let
$\widehat{y}(u)=\widehat{Y}_{\!A}(\ee^{\ii u})$. Then the $m$-th DFT bin of
$\{\widehat{y}(u_k)\}$ equals
\[
\frac{1}{N} \sum_{k=0}^{N-1} \widehat{y}(u_k) \, \ee^{-2\pi\ii mk/N}
=\sum_{\substack{j\ge0\\ j\equiv m \pmod{N}}}
\frac{(-1)^{j}}{j!}\int_0^{A}\tau^{j}y(\tau) \dd\tau.
\]
In particular the bins are (aliased, signed, truncated) \emph{time moments}
of the data, indexed by non-negative integers only.
\end{proposition}

Time moments are the classical carriers of the $s\to0$ expansion of a
transfer function in moment-matching identification and model reduction
\cite{GibilaroLees1969}; Proposition~\ref{prop:moments} shows the circle
pipeline is an FFT reorganization of them. It also shows the pipeline
returns nothing at negative orders---so the $1/s$ of the step input cannot
appear at $\alpha=-1$ and must be handled separately---and nothing at
non-integer orders, foreshadowing Theorem~\ref{thm:aliasing}.

\begin{proposition}[Truncation bias]\label{prop:bias}
Suppose $y(t)=y_\infty+g(t)$ with $|g(t)|\le C\ee^{-\beta t}$, $\beta>0$
(stable step response, $y_\infty=H(0)$). Then
\[
\widehat{Y}_{\!A}(s)\;=\;\frac{y_\infty}{s}\;+\;G(s) \;-\;
\frac{y_\infty \ee^{-sA}}{s} ,
\qquad G(s)=\int_0^{\infty} g(\tau)\ee^{-s\tau}\dd\tau ,
\]
and the entire function $-y_\infty\ee^{-sA}/s + y_\infty/s
= -y_\infty\sum_{m\ge0}\frac{(-A)^{m+1}}{(m+1)!}\, s^{m}$ contributes the
bias
\begin{equation}\label{eq:bias}
b_m \;=\; -y_\infty \, \frac{(-A)^{m+1}}{(m+1)!}
\end{equation}
to the coefficient of $s^{m}$ for every $m\ge0$---hence, up to the
$m+\ell N$ aliasing of Proposition~\ref{prop:moments} (negligible for
$N\gg A$), to DFT bin $m$---while the $g$-tail term is bounded by
$C\, \ee^{-(\beta+\cos u)A}/(\beta+\cos u)$ for $\cos u>-\beta$.
\end{proposition}

For any useful window the bias \eqref{eq:bias} is catastrophic: with
$y_\infty=0.2$ and $A=12$, $b_0=2.4$ against a true coefficient
$a_1=-0.16$ (the naive pipeline's $s$-multiplication shifts $b_m$ up one
bin), and $b_5=-829.4$. The experiment in
Section~\ref{sec:experiments} reproduces \eqref{eq:bias} to all printed
digits. The repair is equally explicit:

\begin{corollary}[Corrected pipeline]\label{cor:corrected}
Estimate $y_\infty$ from the tail of the record, form $g=y-y_\infty$, and
\begin{equation}\label{eq:corrected}
H(\ee^{\ii u_k})\; \approx\; y_\infty \;+\; \ee^{\ii u_k}\,
\widehat{G}_{\!A}(\ee^{\ii u_k}),
\qquad \widehat{G}_{\!A}(s)=\int_0^{A}g(\tau)\ee^{-s\tau}\dd\tau ,
\end{equation}
before the DFT. If every pole of $H$ satisfies
$\operatorname{Re}(\text{pole})<-1$ (more generally, if the circle radius
$r$ is chosen with $r<\beta$), the total error in each recovered
coefficient is $\mathcal{O}(\ee^{-(\beta-1)A})$ plus quadrature error, and
the procedure coincides with Cauchy-integral coefficient extraction
\cite{LynessMoler1967,Henrici1979, Fornberg1981} applied to $H$ itself,
which the data reconstruct as $y_\infty+s\, \widehat{G}_{\!A}(s)$: here
$\widehat{G}_{\!A}(s)$ approximates the transient transform
$G(s)=(H(s)-H(0))/s$, which is analytic at $s=0$, with error bounded by
the $g$-tail term of Proposition~\ref{prop:bias}.
\end{corollary}

The radius condition matters: for slow systems ($\beta\le1$) the truncated
transform of $g$ no longer converges to $G$ as $A$ grows on the arc
$\cos u\le-\beta$, and the evaluation contour must be shrunk to
$|s|=r<\beta$, rescaling the recovered coefficients by $r^{-n}$; the
conditioning trade-offs of the radius choice are analyzed in depth by
Bornemann \cite{Bornemann2011} (see also
\cite{AustinKravanjaTrefethen2014}).

\subsection{The aliasing obstruction}\label{sec:aliasing}

\begin{theorem}[Integer-support aliasing]\label{thm:aliasing}
Let $y\in L^{1}_{\mathrm{loc}}[0, \infty)$ and $A\le\infty$ be such that
$\widehat{y}(u)=\int_0^{A}y(\tau)\,\ee^{-\tau \ee^{\ii u}}\dd\tau$ converges
for all $u$. Then $\widehat{y}(u+2\pi)=\widehat{y}(u)$ identically.
Consequently, if $\widehat{y}(u)=\int_{\R} \mathcal{B}(\alpha)\ee^{\ii
u\alpha}\dd\alpha$ for a tempered distribution $\mathcal{B}$, then
$\operatorname{supp}\mathcal{B}\subseteq\Z$.
\end{theorem}

\begin{proof}
The integrand depends on $u$ only through $\ee^{\ii u}$. A $2\pi$-periodic
tempered distribution has Fourier transform supported on the integer
lattice (its Fourier series), which is the claim.
\end{proof}

\begin{corollary}\label{cor:blind}
No estimator that evaluates the (truncated) Laplace transform of measured
data on the unit circle and inverts in $u$ can produce mass at non-integer
$\alpha$. In particular, fractional orders of a black-box system cannot be
read off as non-integer spectral lines from such a construction, for any
window $A$, sampling rate, or FFT length.
\end{corollary}

The obstruction is a statement about the \emph{readout}, not about the
data: on $\operatorname{Re}s>0$ the Laplace transform determines $y$
uniquely, so the fractional information is present---it is the circle
parametrization that projects it onto $\Z$. What survives the projection is
a decay signature. For $F(s)=s^{-1/2}$, restricting the branch
$\ee^{-\ii u/2}$ to one period $u\in(0,2\pi)$ and expanding in a Fourier
series,
\begin{equation}\label{eq:branchsignature}
c_m=\frac{1}{2\pi}\int_0^{2\pi}\ee^{-\ii u/2}\ee^{-\ii mu} \dd u
=\frac{-\ii}{\pi\, (m+\tfrac12)} ,
\end{equation}
an integer comb decaying only like $1/m$---the fingerprint of a branch cut
crossing the contour, as opposed to the geometric decay of an analytic
$H$. Slow comb decay can thus \emph{detect} non-integer content, but the
order itself cannot be localized this way. When $H$ is known as a formula,
the lift to $u\in\R$ (Section~\ref{sec:representation}) is available and
half-integer orders reappear at period $4\pi$; for data, they do not.

\section{The logarithmic-spiral estimator}\label{sec:spiral}

The failure mechanism of the circle is that $s=\ee^{\ii u}$ revisits the
same point every $2\pi$ while $s^{\alpha}$, $\alpha\notin\Z$, does not.
The fix is to sample where the contour never revisits: a logarithmic
spiral
\begin{equation}\label{eq:spiral}
s_j=\ee^{(\sigma+\ii)u_j}, \qquad u_j=u_0+j\, \Delta u, \quad j=0, \dots,N-1,
\qquad \sigma\neq0 ,
\end{equation}
single-valued along the contour. If the transform under study is a finite
comb with distinct real orders $\alpha_m$, then its samples on the spiral
take the form of a finite exponential sum in the index $j$:
\begin{equation}\label{eq:expsum}
F(s_j) = \sum_{m=1}^{M} c_m\, z_m^{\,j}, \qquad
z_m=\ee^{(\sigma+\ii) \Delta u\, \alpha_m} ,
\end{equation}
an exponential sum in the sample index $j$. Order identification is
thereby converted into \emph{line-spectrum estimation}, for which
classical super-resolution tools exist; we use the matrix-pencil method of
Hua and Sarkar \cite{HuaSarkar1990} (Hankel data matrix, SVD, shifted
singular-subspace pencil), which is exact for noiseless data with
$N\ge2M$ and degrades gracefully with noise.

\begin{proposition}[Identifiability on the spiral]\label{prop:injective}
For $\sigma\neq0$ and distinct real $\alpha_m$, the nodes $z_m$ in
\eqref{eq:expsum} are distinct, and each $\alpha_m$ is uniquely determined
by $z_m$: indeed $\ln|z_m|=\sigma\, \Delta u\, \alpha_m$ is single-valued, so
$\alpha_m=\ln|z_m|/(\sigma\Delta u)$ with no branch condition. If moreover
$|\Delta u\, \alpha_m|<\pi$, the better-conditioned two-component formula
$\alpha_m=\operatorname{Log}z_m/\big((\sigma+\ii) \Delta u\big)$ applies,
with $\operatorname{Log}$ the principal branch.
\end{proposition}

\begin{proof}
$z_m=z_{m'}$ requires $(\sigma+\ii)\Delta u(\alpha_m-\alpha_{m'}) \in
2\pi\ii\Z$; the real part forces $\sigma\Delta u(\alpha_m-\alpha_{m'})=0$,
hence $\alpha_m=\alpha_{m'}$. The moduli identity is immediate, and under
$|\Delta u\, \alpha_m|<\pi$ the exponent lies in the principal strip.
\end{proof}

The proposition isolates the degree of freedom the circle lacks: at
$\sigma=0$ the moduli collapse to $1$ and distinct orders can alias; within
the family $\ee^{(\sigma+\ii)u}$, any $\sigma\neq0$ restores injectivity.
(Sampling on the positive real axis, $s=\ee^{u}$, would also break the
periodicity, at the cost of a classical real-node Prony problem; we
return to a variant of it in Section~\ref{sec:e5}.)

\paragraph{Working from data.} For measured $y$ the values
$F(s_j)=\widehat{Y}_{\!A}(s_j)$ must converge, which requires
$\operatorname{Re}s_j>0$ when $y$ decays slowly (fractional impulse
responses decay algebraically; steady states not at all). We therefore
restrict the spiral to an arc in the right half-plane,
$u\in[u_0, u_0+(N-1) \Delta u]\subset(-\tfrac\pi2, \tfrac\pi2)$, on which
the truncation error is bounded by
$\ee^{-A\min_j \operatorname{Re}s_j}$ up to algebraic factors. Steady states
need no special treatment here: a constant component of $y$ is itself part
of the order comb of $\widehat{Y}$ (mass $y_\infty$ at $\alpha=-1$) and is
recovered as such, since the arc avoids the divergence region that plagued
the circle. Two practical caveats: the model order $M$ must be chosen (the
singular-value profile of the pencil matrix gives a reliable indicator, cf.\
Section~\ref{sec:experiments}), and recovery of \emph{continuous} order
densities from a right-half-plane arc is a Laplace-type inversion and hence
exponentially ill-posed---consistent with the conditional-stability results
for distributed-order inverse problems \cite{LiuSunYamamoto2021}; the
line-spectrum setting (finite combs) is the regime in which the method is
exact.

\section{Numerical experiments}\label{sec:experiments}

All experiments are reproduced by \texttt{order\_spectrum.py} (NumPy only);
E5 by \texttt{cantor\_pencil.py} and \texttt{deaf\_channel.py}. Laplace
transforms of sampled data are computed with composite Simpson
quadrature.

\subsection{E1: integer-order system from its step response}

$H(s)=1/(s^{2}+4s+5)$, poles $-2\pm\ii$ ($\beta=2$,
$|\mathrm{poles}|=\sqrt5$), step response
$y(t)=\tfrac15-\tfrac15\ee^{-2t}(2\sin t+\cos t)$ sampled with
$N_t=65537$ points on $[0,12]$; $N_u=64$ circle points. The steady state is
estimated as the mean of the last $2\%$ of samples. Exact Taylor
coefficients satisfy $5a_n+4a_{n-1}+a_{n-2}=0$ for $n\ge2$, with
$5a_0=1$ and $5a_1+4a_0=0$; the corrected and naive estimates are compared
in Table~\ref{tab:E1} below.

\begin{table}[ht]
\centering\small
\begin{tabular}{rrrr}
\toprule
$n$ & exact $a_n$ & corrected pipeline & naive pipeline \\
\midrule
0 & $+0.2$        & $+0.200000000005$  & $2.5\times10^{-12}$ \\
1 & $-0.16$       & $-0.160000000057$  & $+2.24$ \\
2 & $+0.088$      & $+0.088000000323$  & $-14.31$ \\
3 & $-0.0384$     & $-0.038400001209$  & $+57.56$ \\
4 & $+0.01312$    & $+0.013120003331$  & $-172.8$ \\
5 & $-0.002816$   & $-0.002816007179$  & $+414.7$ \\
6 & $-0.0003712$  & $-0.000371187522$  & $-829.4$ \\
\bottomrule
\end{tabular}
\caption{E1. Corrected pipeline: max error $1.2\times10^{-8}$ with
estimated steady state ($8.1\times10^{-9}$ if $y_\infty$ is known). The
naive pipeline is corrupted exactly as predicted by
Proposition~\ref{prop:bias}: at bin~1 the predicted excess is
$-y_\infty(-A)=+2.400$; the observed excess is $+2.400$. At bin 6,
$b_5=-0.2\cdot12^{6}/6!=-829.44$ matches all printed digits.}
\label{tab:E1}
\end{table}

\subsection{E2: the aliasing signature of a branch cut}

Sampling the branch $F(\ee^{\ii u})=\ee^{-\ii u/2}$ of $s^{-1/2}$ at
$N_u=64$ points of one period and FFT-ing yields bins matching the
closed form \eqref{eq:branchsignature}: $|\mathrm{bin}_0|=0.6367$ vs.\
$1/(\pi\cdot\tfrac12)=0.6366$; $|\mathrm{bin}_3|=0.0914$ vs.\ $0.0909$;
$|\mathrm{bin}_{10}|=0.0317$ vs.\ $0.0303$ (the residuals combine the
$m+\ell N_u$ aliasing sums with a $1/N_u=0.0156$ offset from sampling the
branch-cut discontinuity at $u=0$; the latter dominates in the high bins).
Energy spreads over \emph{all} integer bins with $1/(m+\tfrac12)$ decay;
nothing appears at $\alpha=-\tfrac12$, illustrating
Theorem~\ref{thm:aliasing}.

\subsection{E3: spiral estimator, exact fractional samples}

$H(s)=s^{-1/2}+0.5\,s^{1.3}$ sampled on the spiral $\sigma=0.2$,
$\Delta u=0.3$. With $N=8$ noise-free samples the matrix pencil returns
$\alpha=(-0.500000000000,\,+1.300000000000)$, weights $(1.0, \,0.5)$,
order-position error $5\times10^{-16}$. With $N=64$ and relative sample
noise $10^{-8}$ (resp.\ $10^{-4}$) the order errors are
$4.7\times10^{-9}$ (resp.\ $4.7\times10^{-5}$): errors scale linearly with
the noise floor.

\subsection{E4: fractional identification from raw time samples}

$y(t)=1/\sqrt{\pi t}+2$, whose transform is
$\widehat{Y}(s)=s^{-1/2}+2\,s^{-1}$: a two-line order comb with one
fractional line. The substitution $t=v^{2}$ removes the integrable
singularity ($\int_0^{A}y\,\ee^{-st}\dd t=\int_0^{\sqrt A}(2/\sqrt\pi+4v)\,
\ee^{-sv^{2}}\dd v$), sampled with $N_v=32769$ points, $A=120$. The spiral
arc $\sigma=0.15$, $u\in[-1.35,1.35]$, $N=64$ lies in
$\operatorname{Re}s\ge0.17$; the recovered orders and weights are reported
in Table~\ref{tab:E4} below.

\begin{table}[ht]
\centering\small
\begin{tabular}{ccccc}
\toprule
noise (rel., on $y$) & $\hat\alpha_1$ & $\hat\alpha_2$ &
$\hat c_1$ & $\hat c_2$ \\
\midrule
$0$        & $-1.00000000$ & $-0.50000000$ & $2.00000$ & $1.00000$ \\
$10^{-4}$  & $-0.99997350$ & $-0.49995259$ & $2.00021$ & $0.99979$ \\
\bottomrule
\end{tabular}
\caption{E4. True $(\alpha, c)=(-1,2), (-\tfrac12,1)$. Noise-free order
error $1.8\times10^{-9}$; at $10^{-4}$ noise, $4.7\times10^{-5}$. The
pencil singular values ($97.1, \ 0.89,\ 1.6\times10^{-9}, \dots$
noise-free; third value rising to $7.8\times10^{-6}$ under noise) cleanly
indicate the model order $M=2$.}
\label{tab:E4}
\end{table}

Experiment E4 performs, from time-domain samples alone, precisely the task
that Corollary~\ref{cor:blind} proves impossible for the unit-circle
readout: localization of a non-integer order. The steady-state-like
component ($2/s$) is recovered as the $\alpha=-1$ line rather than
subtracted, as promised in Section~\ref{sec:spiral}.

\subsection{E5: a calibration standard from fractal geometry---complex
dimensions of the Cantor string}\label{sec:e5}

Experiments E1--E4 validate the estimator against synthetic signals drawn
from the very model class the method assumes; the ground truth is planted
by the experimenter. A sharper test requires systems whose exponential-sum
structure is a \emph{theorem} rather than a construction, with parameters
known through an entirely independent computational route, and---ideally---
available in families where a dichotomy predicted by theory (uniform versus
non-uniform damping) is \emph{certified} rather than conjectural. The
theory of fractal strings supplies exactly this: a bounded open set
$\mathcal{L}=\{\ell_j\}\subset\R$ has geometric zeta function
$\zeta_{\mathcal{L}}(s)=\sum_j \ell_j^{s}$---an instance of the
representation \eqref{eq:orderintegral} with order comb at the
log-lengths---whose poles, the \emph{complex dimensions}, govern the
string's geometric and spectral oscillations
\cite{LapidusPomerance1993,LapidusVanFrankenhuijsen2013}.
Self-similar (``lattice'') strings have all complex dimensions on one
vertical line, a theorem rather than a hypothesis; this is the calibration
standard used below.

\paragraph{The system.} The Cantor string (complement of the middle-thirds
Cantor set in $[0,1]$) has lengths $3^{-n}$ with multiplicities $2^{n-1}$,
$\zeta_{\mathcal{L}}(s)=3^{-s}/(1-2\cdot3^{-s})$, and complex dimensions
known in closed form:
\begin{equation}\label{eq:cantordims}
\omega_k = D+\ii kp, \qquad
D=\frac{\log2}{\log3}=0.6309\,2975\ldots,
\end{equation}
\begin{equation*}
p=\frac{2\pi}{\log3}=5.7192\,0173\ldots, \qquad k\in\Z,
\end{equation*}
all on the single vertical line $\operatorname{Re}s=D$: the \emph{lattice
case}. The frequency (spectral) counting function of the associated
Dirichlet Laplacian satisfies the spectral factorization
$\zeta_\nu(s)=\zeta_{\mathcal{L}}(s)\,\zeta(s)$, so each geometric
oscillation reaches the spectrum multiplied by the channel gain
$\zeta(\omega_k)$; the Lapidus--Maier inverse spectral theorem
\cite{LapidusMaier1995} identifies zeros of $\zeta$ with null frequencies
of this geometry-to-spectrum channel, with the critical line as the only
possibly-deaf band.

\paragraph{Method.} In the logarithmic variable $\lambda=\log x$, Poisson
summation gives the Cantor string's geometric measure \emph{exactly} as a
mode expansion,
\begin{equation*}
\sum_{n\ge1}2^{n-1}\delta(\lambda-n\log3)
=\tfrac{1}{2\log3}\sum_{k\in\Z}\ee^{\omega_k\lambda}
\end{equation*}
(no error term---in contrast with arithmetic combs such as the prime-power
comb, where the analogous expansion carries an explicit-formula error
term). Gaussian smoothing by
$\varphi_h$ multiplies mode $k$ by $\ee^{\omega_k^{2}h^{2}/2}$, truncating
the tower to an effective model order exactly as in
Section~\ref{sec:spiral}; sampling at $\lambda_j=\lambda_0+j\Delta$
produces the exponential-sum model \eqref{eq:expsum} with nodes
$z_k=\ee^{\omega_k\Delta}$, and the matrix pencil returns
$\operatorname{Re}\hat\omega=\ln|z|/\Delta$ (the damping/Minkowski-dimension
line) and $\operatorname{Im}\hat\omega=\arg z/\Delta$ (oscillatory
dimensions), by the argument of Proposition~\ref{prop:injective} (with the
roles of contour and orders interchanged: a real logarithmic grid
resolving complex orders). For the spectral run, the
smoothed spectral measure has atoms at $\lambda=\log(m\cdot3^{n})$; its
smooth part is the Weyl term $\ee^{\lambda}$ (total length $1$, the
string's steady state, subtracted exactly as in
Corollary~\ref{cor:corrected}) plus modes weighted by
$\zeta(\omega_k)/(2\log3)$. Parameters throughout: $h=0.10$,
$\Delta=0.07$, $N=64$, $\lambda_0=5.0$, pencil parameter $N/2$,
energy-threshold retention, a priori filter
$\operatorname{Re}\hat\omega\in(0.2,1.2)$.

\paragraph{F1: complex dimensions from geometry.} Input: the lengths only.

\begin{table}[ht]
\centering\small
\begin{tabular}{cccc}
\toprule
$\operatorname{Re}\hat\omega$ & $\operatorname{Im}\hat\omega$ & exact
$\operatorname{Im}=kp$ & $k$ \\
\midrule
$0.630930$ & $0.000000$  & $0.000000$  & $0$ \\
$0.630930$ & $5.719202$  & $5.719202$  & $1$ \\
$0.630930$ & $11.438403$ & $11.438403$ & $2$ \\
$0.630930$ & $17.157605$ & $17.157605$ & $3$ \\
$0.630930$ & $22.876807$ & $22.876807$ & $4$ \\
$0.630930$ & $28.596009$ & $28.596009$ & $5$ \\
$0.630930$ & $34.315211$ & $34.315210$ & $6$ \\
$0.630928$ & $40.034412$ & $40.034412$ & $7$ \\
\bottomrule
\end{tabular}
\caption{E5/F1: complex dimensions of the Cantor string recovered from
geometric data alone. Every real part equals $D=\log2/\log3$ to the
printed six decimals: the lattice theorem measured by the instrument---a
uniformly damped order comb, here certified truth rather than conjecture.}
\label{tab:E5}
\end{table}

Every real part equals $D$ to the printed six decimals: the lattice
theorem, measured by the instrument. A ninth node ($k=8$) emerges degraded
at the smoothing floor $\ee^{-k^{2}p^{2}h^{2}/2}$, exactly where the
weight ladder predicts the retained set to end.

\paragraph{F2: the same dimensions from the sound, with $\zeta$ as the
measured gain.} Input: the spectral counting data, Weyl term subtracted.
The pencil recovers the same tower ($\operatorname{Re}\hat\omega=0.630930$
for $k\le5$, drifting to $0.630937$ at the floor), and the amplitude ratio
of each spectral mode to its geometric counterpart matches the channel
gain independently computed from $\zeta$ to four decimals in every one of
eight modes (e.g.\ $k=0$: measured $2.1598$ vs.\ $|\zeta(D)|=2.1598$;
$k=5$: measured $2.2269$ vs.\ $2.2269$): an \emph{acoustic measurement of
$|\zeta|$ along the vertical line $\operatorname{Re}s=D$}, obtained
entirely from integer counting data, realizing the spectral factorization
$\zeta_\nu=\zeta_{\mathcal{L}}\cdot\zeta$ as a live, measured transfer
function.

\paragraph{F3: the deaf channel at the critical line.} To exhibit the
Lapidus--Maier obstruction directly, a generalized string with geometric
counting function
$N_{\mathcal{L}}(y)=y^{1/2}+\varepsilon\operatorname{Re}(y^{1/2+\ii\gamma})$,
$\varepsilon=0.5$ (geometry oscillating at complex dimension
$\tfrac12+\ii\gamma$), gives, after computing $N_\nu$ by direct summation
(no $\zeta$ anywhere in the forward model) and subtracting the Weyl term: a
control frequency $\gamma=10.0$ (not a zero) survives with measured
amplitude $0.7747$ against a predicted $\varepsilon|\zeta(\tfrac12+10\ii)|
=0.7746$; the geometric oscillation at $\gamma=14.134725$ (the first zeta
zero), of the same input amplitude $0.5$, survives with measured amplitude
$\mathbf{0.0007}$ against a predicted $0.0000$. A fractal string can hide
information from its own sound precisely at the zeta zeros; the
estimator's blind spots, read off as measured channel gains, \emph{are}
the zeros---the inverse-spectral content of the Riemann Hypothesis made
visible as an instrument reading, with no claim of bearing on the
Hypothesis itself: the zeros are the kernel of a forward map, and
information annihilated by a forward map is unrecoverable by any
estimator.

\paragraph{Why this experiment matters for the estimator.} F1--F3 supply
three validation claims that E1--E4 cannot. (i)~\emph{End-to-end
correctness against an external, independently-computed standard}, in a
regime (F1) where the target values are proven rather than conjectured.
(ii)~\emph{The complex-order regime}: the pencil and the modulus formula
of Proposition~\ref{prop:injective} resolve complex orders, and F2 shows
the recovered \emph{amplitudes}, not just positions, carry independently
verifiable physical content (the values of $\zeta$ itself). (iii)~A
\emph{predicted failure mode}, deliberately engineered (F3): the
estimator's channel gain vanishes exactly where theory says it must,
demonstrating that the method's blind spots are structural, not
accidental, and are located precisely where the theory of
Section~\ref{sec:circle}--\ref{sec:spiral} predicts.

\section{Related work}\label{sec:related}

\paragraph{Distributed-order systems and order distributions.}
Equation \eqref{eq:orderintegral} is the transfer function of
distributed-order systems, originating with Caputo
\cite{Caputo1969,Caputo1995,Caputo2001} and developed as an ``order
domain'' by Bagley and Torvik \cite{BagleyTorvik2000a, BagleyTorvik2000b};
Lorenzo and Hartley formalized variable- and distributed-order operators
\cite{LorenzoHartley2002} and the Laplace-domain operator classes that
order distributions generate \cite{HartleyLorenzo2011}, and the monograph
\cite{JiaoChenPodlubny2012} treats stability and simulation of the class;
see \cite{DingReview2021} for a survey. The identification of order
distributions was posed and attacked by Hartley and Lorenzo
\cite{HartleyLorenzo1999,HartleyLorenzo2003}: they already observe the
structural identity behind \eqref{eq:naiveinversion} (the order integral
as a Laplace transform in $-\log s$, their $r$-Laplace transform,
Eq.~(18) of \cite{HartleyLorenzo1999}) and represent integer systems as
impulsive order distributions; their estimator, however, discretizes the
order axis and solves least squares against frequency-response data, as do
its successors \cite{NazarianHaeri2010,DuanLi2021}. Optimization- and
instrumental-variable estimation of discrete fractional orders is
developed in \cite{VictorMalti2013,VictorMalti2022}; log-periodogram
regression estimates long-memory orders in time series
\cite{GewekePorterHudak1983}. The present paper differs in estimator
class: closed-form/spectral rather than fitted, with the obstruction
theorem delimiting where the spectral route is possible.

\paragraph{Cauchy-integral coefficient extraction.} Computing
Taylor--Laurent coefficients by trapezoidal Cauchy integrals on circles is
classical \cite{LynessMoler1967,Henrici1979,Fornberg1981}, with definitive
accuracy and radius-selection analysis in \cite{Bornemann2011} and a
modern roots-of-unity framework in \cite{AustinKravanjaTrefethen2014}.
Corollary~\ref{cor:corrected} identifies our corrected integer-order
pipeline with this method applied to a transform computed from data; the
adjacent inverse direction (Weeks-type inversion, unit-circle FFTs of
Laplace-domain functions) is surveyed in \cite{Weideman1999}, and fully
numerical \emph{forward} Laplace transforms of sampled data are treated in
\cite{WeidemanFornberg2023}. The information content of the integer bins
is that of classical time moments \cite{GibilaroLees1969}.

\paragraph{What is new here.} To our knowledge: the exact bias formula
\eqref{eq:bias} for step-response truncation on circle contours; the
aliasing theorem (Theorem~\ref{thm:aliasing}) with the branch-cut decay
signature \eqref{eq:branchsignature} as the data-side limit of the
unit-circle program; the identifiability statement of
Proposition~\ref{prop:injective}; and the composition
\emph{data $\to$ truncated Laplace transform on a right-half-plane spiral
arc $\to$ matrix-pencil line spectrum $\to$ orders and weights}, validated
from raw time samples; and the use of fractal-string complex dimensions as
a calibration benchmark for order-distribution estimators (E5), including
an acoustic measurement of the Riemann zeta function along a vertical
line and a demonstration of the estimator's channel gain vanishing at a
zeta zero. Each ingredient separately is standard.

\section*{Funding}

This research did not receive any specific grant from funding agencies in
the public, commercial, or not-for-profit sectors.

\section*{Conflict of Interest}

The author declares no conflict of interest.

\section*{Data Availability Statement}

\begin{sloppypar}
No experimental or proprietary data were used in this work; all numerical
experiments (Section~\ref{sec:experiments}) use synthetic signals generated
in closed form by the accompanying code, including E5, whose geometric and
spectral counting data are generated in closed form from the Cantor
string's defining lengths within the code itself. The complete source code
implementing every experiment, together with the exact numerical output
validating it, is openly available at
\url{https://drive.google.com/file/d/1eselh0HetF4mnDnc4iptheXgnBU57oXE/view}
(see Appendix~\ref{app:repro} for details).
\end{sloppypar}

\section*{Declaration of AI-Assisted Technologies}

During preparation of this work the author used a large language model
(LLM) to verify derivations, search prior literature, help derive
Propositions~\ref{prop:moments} and~\ref{prop:bias} and
Theorem~\ref{thm:aliasing}, design and implement the accompanying numerical
code. After using this tool the
author reviewed and edited all content as
needed and takes full responsibility for the content of this publication.

\appendix
\section{Proofs}\label{app:proofs}

\begin{proof}[Proof of Proposition~\ref{prop:moments}]
Expanding $\ee^{-\tau\ee^{\ii u}}=\sum_{j\ge0}\frac{(-\tau)^{j}}{j!}
\ee^{\ii ju}$ (absolutely convergent, uniformly on compacts) and
integrating over $[0,A]$ termwise gives
$\widehat{y}(u)=\sum_{j\ge0}\mu_j\,\ee^{\ii ju}$ with
$\mu_j=\frac{(-1)^{j}}{j!}\int_0^{A}\tau^{j}y(\tau) \dd\tau$; the DFT of a
uniform sampling of $\sum_j\mu_j\ee^{\ii ju}$ aliases $j$ modulo $N$ into
the stated sum.
\end{proof}

\begin{proof}[Proof of Proposition~\ref{prop:bias}]
Write $\int_0^{A}=\int_0^{\infty}-\int_A^{\infty}$ on both components of
$y=y_\infty+g$. The constant part gives
$y_\infty\big(1-\ee^{-sA}\big)/s$, whose expansion
\begin{align*}
\frac{y_\infty(1-\ee^{-sA})}{s}
&=\frac{y_\infty}{s}\Big(1-\sum_{k\ge0}\frac{(-sA)^{k}}{k!} \Big)
=-y_\infty\sum_{k\ge1}\frac{(-A)^{k}s^{k-1}}{k!} \\
&=-y_\infty\sum_{m\ge0}\frac{(-A)^{m+1}}{(m+1)!}\,s^{m}
\end{align*}
is entire: the pole $y_\infty/s$ cancels exactly (which is why no mass
appears at $\alpha=-1$) and the coefficient of $s^{m}$ is
\eqref{eq:bias}. The transient part contributes $G(s)$ minus its tail,
bounded as stated since $|\ee^{-s\tau}|=\ee^{-\tau\cos u}$ on $|s|=1$.
\end{proof}

\emph{Note.} The Poisson-summation and Gaussian-smoothing identities used
in E5 follow the standard explicit-formula machinery for self-similar
fractal strings \cite{LapidusVanFrankenhuijsen2013}: the geometric measure
of the Cantor string is log-periodic after renormalization by
$\ee^{-D\lambda}$, so the Gaussian--exponential integrals evaluate exactly
to
$\int \varphi_h(u)\ee^{-\omega u}\dd u = \ee^{\omega^{2}h^{2}/2}$ for each
complex dimension $\omega=\omega_k$, with no further error term---in
contrast with arithmetic combs such as the prime-power comb.

\section{Reproducibility}\label{app:repro}

\begin{sloppypar}
\texttt{order\_spectrum.py} (NumPy) implements
\texttt{circle\_taylor\_coeffs} (Corollary~\ref{cor:corrected}),
\texttt{naive\_circle\_coeffs} (for Table~\ref{tab:E1}'s bias column),
\texttt{matrix\_pencil} \cite{HuaSarkar1990}, and
\texttt{spiral\_order\_estimate} (Section~\ref{sec:spiral}); running the
file reproduces E1--E4 verbatim. \texttt{cantor\_pencil.py} (NumPy +
mpmath, the latter only for the independent reference $|\zeta|$ values
used to check F2; no zeta values enter the forward data) reproduces F1 and
F2 and Table~\ref{tab:E5}; \texttt{deaf\_channel.py} reproduces F3. Code
available at
\url{https://drive.google.com/file/d/1eselh0HetF4mnDnc4iptheXgnBU57oXE/view}
(\texttt{order\_spectrum.py}, E1--E4),
\url{https://drive.google.com/file/d/1AMGYCC5RAgxeH3odohpvt_nBpiDiR6zQ/view}
(\texttt{cantor\_pencil.py}, F1--F2), and
\url{https://drive.google.com/file/d/1Giu4cQ-A-zB85gdzpZ1xmXw0NMUc6O7U/view}
(\texttt{deaf\_channel.py}, F3).
\end{sloppypar}

\begin{thebibliography}{99}\small

\bibitem{AustinKravanjaTrefethen2014} A.~P.~Austin, P.~Kravanja,
L.~N.~Trefethen, Numerical algorithms based on analytic function values at
roots of unity, \emph{SIAM J.\ Numer.\ Anal.} \textbf{52} (2014)
1795--1821.

\bibitem{BagleyTorvik2000a} R.~L.~Bagley, P.~J.~Torvik, On the existence of
the order domain and the solution of distributed order equations, Part~I,
\emph{Int.\ J.\ Appl.\ Math.} \textbf{2} (2000) 865--882.

\bibitem{BagleyTorvik2000b} R.~L.~Bagley, P.~J.~Torvik, On the existence of
the order domain and the solution of distributed order equations, Part~II,
\emph{Int.\ J.\ Appl.\ Math.} \textbf{2} (2000) 965--988.

\bibitem{Bornemann2011} F.~Bornemann, Accuracy and stability of computing
high-order derivatives of analytic functions by Cauchy integrals,
\emph{Found.\ Comput.\ Math.} \textbf{11} (2011) 1--63.

\bibitem{Caputo1969} M.~Caputo, \emph{Elasticit\`a e Dissipazione},
Zanichelli, Bologna, 1969.

\bibitem{Caputo1995} M.~Caputo, Mean fractional-order-derivatives
differential equations and filters, \emph{Ann.\ Univ.\ Ferrara}
\textbf{41} (1995) 73--84.

\bibitem{Caputo2001} M.~Caputo, Distributed order differential equations
modelling dielectric induction and diffusion, \emph{Fract.\ Calc.\ Appl.\
Anal.} \textbf{4} (2001) 421--442.

\bibitem{DingReview2021} W.~Ding, S.~Patnaik, S.~Sidhardh, F.~Semperlotti,
Applications of distributed-order fractional operators: a review,
\emph{Entropy} \textbf{23} (2021) 110.

\bibitem{DuanLi2021} J.-S.~Duan, Y.~Li, Identification of the system with
distributed-order derivatives, \emph{Fract.\ Calc.\ Appl.\ Anal.}
\textbf{24} (2021) 1619--1628.

\bibitem{Fornberg1981} B.~Fornberg, Numerical differentiation of analytic
functions, \emph{ACM Trans.\ Math.\ Software} \textbf{7} (1981) 512--526.

\bibitem{GewekePorterHudak1983} J.~Geweke, S.~Porter-Hudak, The estimation
and application of long memory time series models, \emph{J.\ Time Ser.\
Anal.} \textbf{4} (1983) 221--238.

\bibitem{GibilaroLees1969} L.~G.~Gibilaro, F.~P.~Lees, The reduction of
complex transfer function models to simple models using the method of
moments, \emph{Chem.\ Eng.\ Sci.} \textbf{24} (1969) 85--93.

\bibitem{HartleyLorenzo1999} T.~T.~Hartley, C.~F.~Lorenzo, \emph{Fractional
System Identification: An Approach Using Continuous Order-Distributions},
NASA/TM-1999-209640, NASA Glenn Research Center, 1999.

\bibitem{HartleyLorenzo2003} T.~T.~Hartley, C.~F.~Lorenzo, Fractional-order
system identification based on continuous order-distributions,
\emph{Signal Process.} \textbf{83} (2003) 2287--2300.

\bibitem{HartleyLorenzo2011} T.~T.~Hartley, C.~F.~Lorenzo,
Order-distributions and the Laplace-domain logarithmic operator,
\emph{Adv.\ Difference Equ.} \textbf{2011}:59 (2011).

\bibitem{Henrici1979} P.~Henrici, Fast Fourier methods in computational
complex analysis, \emph{SIAM Rev.} \textbf{21} (1979) 481--527.

\bibitem{HuaSarkar1990} Y.~Hua, T.~K.~Sarkar, Matrix pencil method for
estimating parameters of exponentially damped/undamped sinusoids in
noise, \emph{IEEE Trans.\ Acoust.\ Speech Signal Process.} \textbf{38}
(1990) 814--824.

\bibitem{JiaoChenPodlubny2012} Z.~Jiao, Y.~Chen, I.~Podlubny,
\emph{Distributed-Order Dynamic Systems: Stability, Simulation,
Applications and Perspectives}, SpringerBriefs in Control, Automation and
Robotics, Springer, London, 2012.

\bibitem{LapidusMaier1995} M.~L.~Lapidus, H.~Maier, The Riemann hypothesis
and inverse spectral problems for fractal strings, \emph{J.\ London Math.\
Soc.} (2) \textbf{52} (1995) 15--34.

\bibitem{LapidusPomerance1993} M.~L.~Lapidus, C.~Pomerance, The Riemann
zeta-function and the one-dimensional Weyl--Berry conjecture for fractal
drums, \emph{Proc.\ London Math.\ Soc.} (3) \textbf{66} (1993) 41--69.

\bibitem{LapidusVanFrankenhuijsen2013} M.~L.~Lapidus, M.~van
Frankenhuijsen, \emph{Fractal Geometry, Complex Dimensions and Zeta
Functions: Geometry and Spectra of Fractal Strings}, 2nd Ed., Springer
Monographs in Mathematics, Springer, New York, 2013.

\bibitem{LiuSunYamamoto2021} J.~Liu, C.~Sun, M.~Yamamoto, Recovering the
weight function in distributed order fractional equation from interior
measurement, \emph{Appl.\ Numer.\ Math.} (2021).

\bibitem{LorenzoHartley2002} C.~F.~Lorenzo, T.~T.~Hartley, Variable order
and distributed order fractional operators, \emph{Nonlinear Dynam.}
\textbf{29} (2002) 57--98.

\bibitem{LuchkoKiryakova2013} Y.~Luchko, V.~Kiryakova, The Mellin integral
transform in fractional calculus, \emph{Fract.\ Calc.\ Appl.\ Anal.}
\textbf{16} (2013) 405--430.

\bibitem{LynessMoler1967} J.~N.~Lyness, C.~B.~Moler, Numerical
differentiation of analytic functions, \emph{SIAM J.\ Numer.\ Anal.}
\textbf{4} (1967) 202--210.

\bibitem{NazarianHaeri2010} P.~Nazarian, M.~Haeri, Generalization of order
distribution concept used in the fractional order system identification,
\emph{Signal Process.} \textbf{90} (2010) 2243--2252.

\bibitem{VictorMalti2013} S.~Victor, R.~Malti, H.~Garnier, A.~Oustaloup,
Parameter and differentiation order estimation in fractional models,
\emph{Automatica} \textbf{49} (2013) 926--935.

\bibitem{VictorMalti2022} S.~Victor, A.~Mayoufi, R.~Malti, M.~Chetoui,
M.~Aoun, System identification of MISO fractional systems: parameter and
differentiation order estimation, \emph{Automatica} \textbf{141} (2022)
110268.

\bibitem{Weideman1999} J.~A.~C.~Weideman, Algorithms for parameter
selection in the Weeks method for inverting the Laplace transform,
\emph{SIAM J.\ Sci.\ Comput.} \textbf{21} (1999) 111--128.

\bibitem{WeidemanFornberg2023} J.~A.~C.~Weideman, B.~Fornberg, Fully
numerical Laplace transform methods, \emph{Numer.\ Algorithms}
\textbf{92} (2023) 985--1006.

\end{thebibliography}
\end{document}
